\section{Problem 218: prime-scale growth does not turn oscillation into density}
\subsection*{1) OBJECTIVE AND SOURCE REVIEW}
Attempt dated 2026-09-23. For consecutive prime gaps $d_n$, the target concerns density $1/2$ for each of $d_{n+1}\ge d_n$ and $d_{n+1}\le d_n$. I read \texttt{gpt\_pro\_5.4/218.tex} completely. It invokes Pintz to obtain arbitrarily large and small consecutive-gap ratios. I test whether those oscillations together with the first-order total gap scale can imply the requested frequencies.
\subsection*{2) WORK}
Here is an explicit obstruction to that inference. Define positive even model gaps by
\[
g_n=2\lceil\log(n+2)\rceil\quad\text{unless }n=2^j\ (j\ge2),
\qquad g_{2^j}=2j^4.
\]
The baseline has total $2N\log N+O(N)$, by comparison of $\sum\log(n+2)$ with its integral and bounded rounding errors. Modifications at powers of two contribute $O((\log N)^5)$ up to $N$. Hence
\[
\sum_{n\le N}g_n\sim2N\log N.
\]
The baseline changes only $O(\log N)$ times up to $N$, and the modifications affect only $O(\log N)$ consecutive comparisons. Thus
\[
\#\{n\le N:g_{n+1}\ne g_n\}=O(\log N).
\]
Both weak-comparison sets therefore have density 1, while both strict-comparison sets have density zero. Nevertheless, immediately before $2^j$ the ratio is of order $j^3$, and immediately after $2^j$ it is of order $j^{-3}$. Therefore
\[
\limsup g_{n+1}/g_n=\infty,\qquad\liminf g_{n+1}/g_n=0.
\]
The model has positivity, evenness, prime-scale total growth, and frequent-scale extreme ratios, but completely different comparison densities. These elementary properties cannot settle the prime question.
\subsection*{3) ADVERSARIAL VERIFICATION}
This is explicitly a model of gaps, not a sequence of actual prime gaps or a counterexample to the conjecture. The total-growth coefficient is immaterial to the logical obstruction. The equality set matters: for any sequence the two weak-comparison indicators add to one plus the equality indicator. Accordingly their two proposed densities of $1/2$ require the equality set to have density zero, a fact not supplied by oscillation alone.
\subsection*{4) FINAL}
\textbf{UNRESOLVED.} Actual prime-specific frequency estimates, including control of equal consecutive gaps, are still missing. The additional result is an explicit limitation on what the supplied qualitative conclusions can imply.
\subsection*{5) SOURCES CHECKED}
Complete supplied file; J. Pintz, \emph{On the ratio of consecutive gaps between primes}, \texttt{https://arxiv.org/abs/1406.2658}, checked 2026-09-23 for the qualitative extreme-ratio theorem. The model construction is proved here.
\subsection*{6) COMPLETION ESTIMATE}
Subjective progress toward the density assertions.
\textbf{COMPLETION: 10\%}
